\section{Brownian observations and their innovations}
\label{sec:observation}

Subtracting the posterior mean drift from the observations will produce
a family of independent Brownian paths. Section~\ref{sec:decoder} will
use a uniform response bound to reconstruct the observations
deterministically from this family.

The finite-dimensional innovation identity is classical in nonlinear
filtering; see Fujisaki, Kallianpur, and
Kunita~\cite[Lemma~2.2]{FujisakiKallianpurKunita1972}.
Here the probabilistic argument applies to any probability law
$\nu$ on $\{0,1\}^I$, where $I$ is countable. For each $i\in I$, fix finite
sets $S_n(i)$ increasing to $I$. Define $b_{i,S}$ by the finite-tilt formula
\eqref{eq:finite-tilt}, with this law $\nu$, and define $b_i$ by
\eqref{eq:drift}. These functions are bounded and jointly Borel. No response
bound is needed in this section.

On a reference probability space let $X\sim\nu$, and independently let
$(B_i)_{i\in I}$ be independent standard Brownian motions. Set
\begin{equation}\label{eq:observation}
 Y_i(t)=tX_i+B_i(t),\qquad
 \mathcal H_t=\sigma\{Y_i(s):i\in I,\ 0\le s\le t\}.
\end{equation}
We use these raw observation sigma-fields; no completion or right-continuity
assumption on the filtration is necessary.

\begin{lemma}[Posterior drift]\label{lem:posterior}
For every fixed $t\ge0$ and $i\in I$,
\begin{equation}\label{eq:posterior}
 b_i(t,Y(t))=\E[X_i\mid\mathcal H_t]\qquad\text{almost surely}.
\end{equation}
\end{lemma}

\begin{proof}
For finite $S\subset I$, write
$\mathcal H_t^S=\sigma\{Y_j(s):j\in S,\ 0\le s\le t\}$.
When $t>0$, the density of the endpoint vector $(Y_j(t))_{j\in S}$,
conditional on $X$, relative to independent $N(0,t)$ variables is
\[
 \exp\!\left(\sum_{j\in S}(y_j-t/2)X_j\right).
\]
Here $X_j^2=X_j$. Moreover,
\[
 Y_j(s)-\frac{s}{t}Y_j(t)
 =B_j(s)-\frac{s}{t}B_j(t),\qquad 0\le s\le t.
\]
These bridge paths are jointly independent of $X$ and the endpoint vector.
Indeed, at finitely many times this follows from zero Gaussian covariances
and independence of $B$ from $X$; continuity and rational-time evaluation
extend it to the paths. Bayes' formula therefore gives
\[
 b_{i,S}(t,Y(t))=\E[X_i\mid\mathcal H_t^S]
 \qquad\text{almost surely}.
\]
At $t=0$ both sides equal $\E X_i$.

For fixed $i,t$, the sigma-fields $\mathcal H_t^{S_n(i)}$ increase to generate
$\mathcal H_t$. The bounded conditional expectations
$M_n=\E[X_i\mid\mathcal H_t^{S_n(i)}]$ converge almost surely to
$\E[X_i\mid\mathcal H_t]$. We recall a short verification of this last step.
For $m\ge n$,
\[
 \E|M_m-M_n|^2=\E M_m^2-\E M_n^2.
\]
Choose $n_k\uparrow\infty$ so that this difference is at most $2^{-4k}$
whenever $m\ge n=n_k$. The maximal inequality obtained by conditioning on
the first crossing gives
\[
 \P\!\left(\max_{n_k\le r\le m}|M_r-M_{n_k}|>2^{-k}\right)
 \le 2^{2k}\E|M_m-M_{n_k}|^2\le2^{-2k}.
\]
To see the first inequality, split into the disjoint events of a first
crossing at $r$ and apply conditional Jensen to $M_m-M_{n_k}$ at that stage.
Letting $m\to\infty$ and summing over $k$ proves that the full sequence is
almost surely Cauchy. Its bounded limit has the same integral as $X_i$
against every event in each finite-stage sigma-field. A monotone-class
argument identifies it with $\E[X_i\mid\mathcal H_t]$.
Thus the limsup defining $b_i$ is an actual limit along the observations
at this fixed time, and \eqref{eq:posterior} follows.
\end{proof}

Define the \emph{innovations} by ordinary scalar time integrals:
\begin{equation}\label{eq:innovation}
 \widehat W_i(t)=Y_i(t)-\int_0^t b_i(s,Y(s))\,ds.
\end{equation}

\begin{lemma}[Joint innovation law]\label{lem:innovation}
The family $(\widehat W_i)_{i\in I}$ has the law of independent standard
Brownian motions. Its finite-coordinate increments after time $u$ are
independent of $\mathcal H_u$.
\end{lemma}

\begin{proof}
Fix $T<\infty$. The map $(s,\omega)\mapsto Y(s,\omega)$ on $[0,T]$ is
measurable for $\mathcal B([0,T])\otimes\mathcal H_T$, coordinate by
coordinate. Composing with the bounded Borel drift shows that the integral
in \eqref{eq:innovation} is defined, measurable from the observation history,
and continuous in its upper endpoint. Hence $\widehat W$ is coordinatewise
continuous and adapted.

For $0\le u<v$ and any bounded $\mathcal H_u$-measurable random variable
$K$, Lemma~\ref{lem:posterior} gives, separately for each $s\in[u,v]$,
\[
 \E[Kb_i(s,Y(s))]=\E[KX_i].
\]
Bounded Fubini integrates these numerical identities. Since
$B_i(v)-B_i(u)$ is independent of $X$ and the Brownian histories through
$u$, it is independent of $\mathcal H_u$. Consequently
\begin{equation}\label{eq:innovation-mean}
 \E[\widehat W_i(v)-\widehat W_i(u)\mid\mathcal H_u]=0.
\end{equation}
This argument uses no single exceptional set valid for every real time.

We identify the joint law using conditional characteristic functions: the
error on each small interval is of order $h^{3/2}$, so the accumulated error
vanishes under subdivision. Let $\theta=(\theta_i)_{i\in I}\in\R^I$ have finite
support, and put
\[
 L_\theta=\sum_i|\theta_i|,\qquad
 \sigma_\theta^2=\sum_i\theta_i^2,\qquad h=v-u.
\]
The increment
\[
 D_{u,v}=\sum_i\theta_i(\widehat W_i(v)-\widehat W_i(u))=Q+R
\]
has a conditionally Gaussian term
$Q=\sum_i\theta_i(B_i(v)-B_i(u))$ of conditional variance
$h\sigma_\theta^2$, and a correction satisfying $|R|\le hL_\theta$.
It follows, without assuming that $Q$ and $R$ are independent, that
\begin{align*}
 \left|\E[D_{u,v}^2\mid\mathcal H_u]-h\sigma_\theta^2\right|
 &\le 2L_\theta\sigma_\theta h^{3/2}+L_\theta^2h^2,\\
 \E[|D_{u,v}|^3\mid\mathcal H_u]
 &\le 4\bigl(c_3\sigma_\theta^3h^{3/2}+L_\theta^3h^3\bigr),
\end{align*}
where $c_3=\E|N(0,1)|^3$. Together with
\eqref{eq:innovation-mean}, the Taylor bound
$|e^{\mathrm ix}-1-\mathrm ix+x^2/2|\le |x|^3/6$ yields a deterministic
constant $C_\theta$ such that, for $h\le1$,
\begin{equation}\label{eq:innovation-characteristic}
 \left|\E[e^{\mathrm iD_{u,v}}\mid\mathcal H_u]
       -e^{-h\sigma_\theta^2/2}\right|
 \le C_\theta h^{3/2}.
\end{equation}

Fix $s<t$ and $A\in\mathcal H_s$. Partition $[s,t]$ into $N$ equal intervals
with endpoints $t_0=s,\ldots,t_N=t$ and length $h\le1$. Set
\[
 a=e^{-h\sigma_\theta^2/2},\qquad
 \mu_k=\E\!\left[\1_A
        \prod_{r=1}^k e^{\mathrm iD_{t_{r-1},t_r}}\right].
\]
The preceding factors are $\mathcal H_{t_{k-1}}$-measurable and have modulus
one. Thus \eqref{eq:innovation-characteristic} gives
$\mu_k=a\mu_{k-1}+\delta_k$ with
$|\delta_k|\le C_\theta h^{3/2}$. Iterating, and using $0<a\le1$, gives
\[
 \left|\E[\1_Ae^{\mathrm iD_{s,t}}]
       -\P(A)e^{-(t-s)\sigma_\theta^2/2}\right|
 \le NC_\theta h^{3/2}\longrightarrow0.
\]
For every finite set of coordinates and every $A\in\mathcal H_s$, uniqueness
of finite-dimensional characteristic functions identifies the increment law
restricted to $A$ as $\P(A)$ times the centered Gaussian law with covariance
$(t-s)\operatorname{Id}$. This proves the vector increment law and independence
from the entire observation past. Repeating over successive time intervals
gives all joint finite-dimensional Brownian laws. Continuous paths, starting
at zero, and rational-time evaluations determine the law on the countable
product of continuous-path spaces. This proves the claim.
\end{proof}
